%%
%% Meta: TI nSpire Einführung
%%       Ziel: Damit die Grundoperationen damit durchgeführt werden können.
%%             Damit man sich an den Rechner gewöhnt.
%%

\input{bmsLayoutPage}

%%%%%%%%%%%%%%%%%%%%%%%%%%%%%%%%%%%%%%%%%%%%%%%%%%%%%%%%%%%%%%%%%%
\renewcommand{\metaHeaderLine}{Arbeitsblatt}
\renewcommand{\arbeitsblattTitel}{Höhenmessung}

\begin{document}%%
\arbeitsblattHeader{}


%% Bild

\begin{center}
\raisebox{1cm}{\includegraphics[width=8cm]{img/Hoehenmessung.png}}
\end{center}

\newcommand{\noteSpace}{\noTRAINER{\rule{0pt}{6.5ex}}}

\begin{tabular}{|c|c|c|c|c|c|c|}
\hline
Erfasser & Distanz $d$ & Augenhöhe & Elevation & Gegenkathete & Totalhöhe $H$ \\
\hline
Name(n) & $d$ & $a$ & $\varepsilon$ &  $g=\LoesungsRaum{d\cdot tan(\epsilon)}$ & $H=\LoesungsRaum{a + g}$ \\
\hline
\noteSpace{}&&&&&\\
\hline
\noteSpace{}&&&&&\\
\hline
\noteSpace{}&&&&&\\
\hline
\noteSpace{}&&&&&\\
\hline
\noteSpace{}&&&&&\\
\hline
\noteSpace{}&&&&&\\
\hline
\noteSpace{}&&&&&\\
\hline
\noteSpace{}&&&&&\\
\hline
\noteSpace{}&&&&&\\
\hline
\end{tabular}

\end{document}
